\paragraph{Matrix gauge transformations and the frame index in the conjugate-coordinate class (conditional; bounded obstruction).}
Work in the Poisson limit on the leaf \(\mathbb{CP}^{1,2}\) with \(T^a\), \(y^\mu\), \(m=M^{45}=y^4\), the sign \(s\), \(\ell^2=-y_ay^a\) and \(R=1\) of the paragraphs on the conjugate-coordinate class in Section~\ref{sec:status-F}, so that \(y\cdot y:=y_\mu y^\mu=-\ell^2-m^2\). With \(\sigma=\pm1\) an orientation fixed by the realisation and \(\epsilon^{0123}=1\),
\[
M^{\mu\nu}=\frac{m\,(y^\mu T^\nu-y^\nu T^\mu)+\sigma\ell\,\epsilon^{\mu\nu\rho\lambda}y_\rho T_\lambda}{y\cdot y},
\]
by hand at \(y^i=0\) and on the explicit realisation at random points (residual \(10^{-12}\)). A first-harmonic gauge parameter is \(\Lambda=\xi_\mu(y)T^\mu\), defined modulo \(\xi_\mu\to\xi_\mu+k\,y_\mu\); take \(\xi\cdot y=0\). The Lie--Poisson relations give exactly
\[
\{\Lambda,T^\nu\}=\lambda^\nu{}_\sigma T^\sigma-s\,m\,T^\sigma\partial^\nu\xi_\sigma,\qquad
\lambda^{\nu\sigma}=\frac{s}{y\cdot y}\bigl[m(\xi^\nu y^\sigma-y^\nu\xi^\sigma)+\sigma\ell\,\epsilon^{\mu\nu\rho\sigma}\xi_\mu y_\rho\bigr],
\]
\[
\{\Lambda,y^\mu\}=s\,m\,\xi^\mu-s\,T^\sigma\partial_\rho\xi_\sigma M^{\rho\mu}.
\]
The first term of the first equation is the algebraic part left uncomputed in the paragraph on the first-order response. It is a local rotation of the index of \(T\), by a bivector \(\lambda(\xi)\) of the three-dimensional space \(W_y\), with three independent functions. It is not independent of the displacement of the base point: \(y_\nu\lambda^{\nu\sigma}=-s\,m\,\xi^\sigma\), so the boost part of \(\lambda\) is what keeps \(T\) orthogonal to the displaced \(y\). At \(y^i=0\) one unit of \(\xi_i\) gives a boost \(m/y^0\) and a rotation \(\ell/y^0\). Also \(W_y\) is not a subalgebra of \(\mathfrak{so}(3,1)\) for \(m\neq0\): for \(aK(\xi)+bJ(\xi)\), with \(K,J\) the boosts and rotations about \(y\), the commutator is \((b^2-a^2)J(\xi\times\xi')+2abK(\xi\times\xi')\) up to a common sign, which stays in the span only for \(a=0\) (exact span of dimension six at \((y^0,y^4)=(5,3),(13,5),(5,4),(25,7)\)). For a plane wave, the derivative term exceeds the rotation by the factor \(\sinh\tau\cosh\tau\,|q|\ell\) with \(\sinh\tau=m/\ell\), so in the regime \(|q|\ell\gg1\) of the frozen symbol the transformation is the Maxwell shift of the paragraph on the tree-level form, and the rotation is of relative order \(1/(|q|\ell)\) there.

Fibre orders do not mix in this action. For \(\zeta(y)\) with \(\zeta\cdot y=\zeta\cdot\zeta=0\), the function \((\zeta\cdot T)^J\) is of fibre order \(J\), and \(\{f(\zeta\cdot T)^J,T^a\}\) is of fibre order \(J\): the terms produced are \((\zeta\cdot T)^{J-1}(\eta\cdot T)\) with \(\eta\cdot y=\eta\cdot\zeta=0\), and their Laplacian on the fibre vanishes, using \(\zeta^\mu\partial\zeta_\mu=0\) for the term with a derivative of \(\zeta\) (numerical projection on the fibre for \(J\le4\) agrees). So at linear order about \(T^a\) only \(J=1\) acts on the frame amplitudes. Moreover no nonconstant gauge parameter fixes the base point: \(\{\Lambda,y^\mu\}=0\) implies \(\{\Lambda,M^{\mu\nu}\}=0\) by the Jacobi identity, hence \(\{\Lambda,T^\mu\}=0\) from \(T^\mu=-M^{\mu\sigma}y_\sigma/y^4\), and \(y^\mu,M^{\mu\nu},T^\mu\) with \(y^4\), a function of \(y^\mu\) on the sheet, separate the points of the sheet. An \(\mathrm{SO}(4,1)\)-equivariant base on which a copy of \(\mathfrak{so}(3,1)\) acts vertically would be a map of the sheet onto \(\mathrm{SO}(4,1)/\mathrm{SO}(3,1)\); none exists, since the stabiliser \(U(2)\) of a point of the sheet has a compact torus of rank two and the compact subgroups of \(\mathrm{SO}(3,1)\) have rank one. Within the Poisson limit, this single orbit, the base functions \(y^\mu\) and linear order about \(T^a\), the matrix gauge group therefore gives no six-parameter local rotation of the frame index. The hypothesis this rests on is the representation.

\paragraph{First-order equivariant class: local \(\mathrm{SO}(3,1)\) inside the unitary group (conditional; not closure).}
Keep Equation~\eqref{eq:Smatrix} with \(C_{\mu\nu\lambda}=0\) and, except where a fermion is named, without the matter term, and change only the Hermitian operators. Let \(\mathcal A\) be a commutative algebra of operators on the Hilbert space of Definition~\ref{def:matrix-coords}, with \([\hat X_a,\mathcal A]\subset\mathcal A\), let \(\mathcal S_{cd}\) be anti-Hermitian derivations of \(\mathcal A\) forming a free and proper unitary action of \(\mathrm{SO}(3,1)\), and write \(\hat X_a=\mathrm i\mathbb D_a\), with
\[
[\mathcal S_{cd},\mathbb D_a]=\eta_{ad}\mathbb D_c-\eta_{ac}\mathbb D_d,\qquad
[\mathbb D_a,\mathbb D_b]=-\mathsf T_{ab}{}^c\mathbb D_c+\tfrac12\mathsf R_{ab}{}^{cd}\mathcal S_{cd}+\Phi_{ab},
\]
with coefficients in \(\mathcal A\), and the \(\mathbb D_a\), \(\mathcal S_{cd}\) pointwise independent derivations on the spectrum \(Z\) of \(\mathcal A\). This absolute parallelism makes \(Z\) a principal bundle over \(Z/\mathrm{SO}(3,1)\) carrying a frame and a connection, by the standard reconstruction of a Cartan connection. These are outputs of the operators, and the action contains no metric. In coordinates, \(\mathbb D_a=u_a{}^be_b{}^\mu(\partial_\mu-u_f{}^g\omega_\mu{}^c{}_g\,\partial/\partial u_f{}^c)+\mathsf t_a-\mathrm iA_a\) with \(u\) a Lorentz matrix, \(\mathsf t_a=\tfrac12\mathsf T_{ac}{}^c\) from Hermiticity (the divergence of the first term is \(\mathsf T_{ac}{}^c\), checked exactly) and \(A_a\) a real order-zero field. This is the construction of~\cite{HanadaKawaiKimura2006}, there on \(\mathrm{Spin}(10)\) with the Euclidean form, read here with the Lorentzian contraction of the matrix index; the boost generators are anti-Hermitian on this Hilbert space. The kinetic tensor \(K^{\mu\nu}=\eta^{ab}\mathbb D_ay^\mu\mathbb D_by^\nu=\eta^{cd}e_c{}^\mu e_d{}^\nu\) does not depend on the fibre, and its Lorentzian signature comes from \(\eta\), so Assumption~\ref{ass:lorentzian} is retained. The merged obstruction for matrices as base coordinates does not apply, since \(y^\mu\) is not among the \(\hat X_a\).

Local rotations are then matrix gauge transformations. For \(\lambda^{cd}=u^c{}_fu^d{}_g\hat\lambda^{fg}(y)\) with six arbitrary functions \(\hat\lambda\), \([\tfrac12\lambda^{cd}\mathcal S_{cd},\mathbb D_a]=\lambda^c{}_a\mathbb D_c-\tfrac12(\mathbb D_a\lambda^{cd})\mathcal S_{cd}\), which rotates \(e\) and shifts \(\omega\) by \(\nabla\hat\lambda\); the finite form is the pull-back by a vertical automorphism of \(Z\), a unitary operator, so \((e,\omega)\) related by a local rotation are one gauge orbit. The coefficient \(\lambda^c{}_a\) lies in the commutative \(\mathcal A\), so the Leibniz defect of the central Moyal branch in the opening paragraphs of this section is absent. The fibre coordinate \(u\) has the index structure of the Stueckelberg field of Equation~\eqref{eq:dressed-omega}, here as an element of the operator algebra; this is an observation and not an identification with Presentation~B. Longitudinal diffeomorphisms are also present, through \(\Lambda=\tfrac12\{\xi^\mu,-\mathrm i\partial_\mu\}\). The field content of the class is \(e\) (16), \(\omega\) (24), \(A\) (4), with gauge parameters \(6+4+1\), and the background-gauge ghost operator on End of the Hilbert space is \([\hat X^a,[\hat X_a,\cdot\,]]\). On \(\Psi=\rho(u)\psi(y)\), with \(\gamma^au_a{}^b=\rho(u)\gamma^b\rho(u)^{-1}\), one has \(\gamma^a\mathbb D_a\Psi=\rho(u)\gamma^be_b{}^\mu(\partial_\mu+\tfrac14\omega_{\mu cd}\gamma^c\gamma^d)\psi\) at \(\mathsf t=A=0\), the Dirac operator with this connection, and an element of \(\mathcal A\) in any representation of \(\mathrm{SO}(3,1)\) sees \(\nabla^a\nabla_a\). A constant \(c\)-number \(\Theta_{ab}\) as in Equation~\eqref{eq:NC} is also available on the flat bundle, through \(A_a\propto\Theta_{ab}(u^{-1}y)^b\), which leaves \(K\) unchanged and breaks the global diagonal Lorentz group to the stabiliser of \(\Theta\); curved backgrounds have operator-valued commutators, and the scale of Equation~\eqref{eq:LUV-def} is then not defined by them.

\emph{Matrix equation.} Since the trace of the action diverges on these operators, as stated in~\cite{HanadaKawaiKimura2006}, read the field equation as stationarity against variations of finite rank. From the two relations alone, \([\mathbb D^a,[\mathbb D_a,\mathbb D_b]]\) equals
\[
\bigl(-\mathbb D^a\mathsf T_{ab}{}^e+\mathsf T_{ab}{}^c\mathsf T^a{}_c{}^e+\mathsf R_{ab}{}^{ae}\bigr)\mathbb D_e+\tfrac12\bigl(\mathbb D^a\mathsf R_{ab}{}^{de}-\mathsf T_{ab}{}^c\mathsf R^a{}_c{}^{de}\bigr)\mathcal S_{de}+\bigl(\mathbb D^a\Phi_{ab}-\mathsf T_{ab}{}^c\Phi^a{}_c\bigr),
\]
each part checked at random frames and connections unrelated to each other (exact on jets; the first two in dimensions three and four, the third in dimension three). Since \(\mathbb D_e\) and \(\mathcal S_{de}\) are independent, each bracket vanishes separately. For \(\mathsf T=0\) the first is the Ricci tensor, \(\mathsf R_{ab}{}^{ae}=-\mathrm{Ric}_b{}^e\), the second follows from it by the Bianchi identity, and the third is \(\nabla^a\Phi_{ab}=0\): in the torsion-free sector the matrix equation is \(\mathrm{Ric}=0\) with a free Maxwell field, the result of~\cite{HanadaKawaiKimura2006}. This was checked as an operator identity on the bundle in exact arithmetic for Ricci-flat pp-waves, Schwarzschild and Kasner, and for non-Ricci-flat pp-waves, \(S^2\times\mathbb R^2\) and de Sitter space, where \([\mathbb D^a,[\mathbb D_a,\mathbb D_b]]=-\mathrm{Ric}_b{}^e\mathbb D_e\) with the Ricci tensor computed separately. No bilinear term is needed for these solutions, and the equation holds with its second-order form, so the order count of the paragraph on the tree-level equations of the level-one amplitudes, whose hypotheses include the orbit, does not arise. With the bilinear term of the conjugate-coordinate class, \(\mathrm{Ric}_{ab}=-\mu^2\eta_{ab}\). An example links the two classes: \(T_a=E_{a4}\), \(M_{ab}=E_{ab}\) in the regular representation of \(\mathrm{SO}(4,1)\) obey the same \(\{T_a,\{T^a,T^b\}\}=-3T^b\) with the same \(\mu^2=-3\), now with base functions \(Y^A=g^A{}_4\), annihilated by \(M_{ab}\), and \(\eta^{ab}g^A{}_ag^B{}_b=\eta^{AB}-Y^AY^B\): unit de Sitter space, with \(\mathrm{Ric}_{ab}=3\eta_{ab}\). How the minimal representation sits in the regular one is not examined. On a background where the bivector map of \(\mathsf R\) is invertible, the generators are themselves built from commutators, \(\mathcal S_{cd}\propto(\mathsf R^{-1})_{cd}{}^{ab}[\mathbb D_a,\mathbb D_b]\) at \(\mathsf T=\Phi=0\): rank six on the sampled points of Schwarzschild and Kasner, rank two on the Ricci-flat pp-wave and zero on flat space.

\emph{Spectrum about flat space (exact, linearised, without a norm).} With \(\mathbb D_a\to\mathbb D_a-h_a{}^c\mathbb D_c+\tfrac12\omega_a{}^{cd}\mathcal S_{cd}\), for \(q^2\neq0\) the kernel is the ten gauge directions (rank 30 of 40, at sampled \(q\)). At null \(q\) there are fourteen further solutions. Two have \(\mathsf T=0\) and are the transverse-traceless polarisations, simple zeros at \(q^2=0\). Twelve have torsion, one of them with totally antisymmetric torsion, which agrees in type with the propagating rank-three field noted in~\cite{HanadaKawaiKimura2006}. Six of the twelve lie at the head of a chain of length two along two transversal lines through the null \(q\), with no longer chain: there the operator has a double zero at \(q^2=0\), with solutions linear in \(y\) times the wave. These six heads have \(\mathsf R=0\) and \(\mathsf T\neq0\). So the set \(\mathsf T=0\) solves the equations and carries only the two simple zeros, while the first-order class is not torsion-free as a set of solutions. With no action defined on the class, no sign of a norm follows for these modes, and this is not a statement about an Ostrogradsky ghost in the sense of Section~\ref{sec:status-N}.

\emph{Sources.} Fields of order zero do not source the Ricci tensor at tree level: for \(\hat X^i=\phi^i\otimes L^i\) with \(\phi^i\in\mathcal A\) and fixed Hermitian \(L^i\), the term \([\hat X^i,[\hat X_i,\hat X_b]]\) is of order zero, and \(\Phi_{ab}\) enters only the order-zero part, so \(\mathrm{Ric}=0\) stands, and likewise for fermion bilinears with \(\Psi\) in \(\mathcal A\). This agrees with the statement of~\cite{HanadaKawaiKimura2006} that the classical equation has no energy-momentum term. It is not a general obstruction: with further matrices that are first-order operators along a compact factor, the same identity gives \(\mathrm{Ric}^{(D)}=0\), whose reduction carries the stress of the Kaluza--Klein fields. For five matrices, \(\mathrm ds^2=\mathrm dU(\mathrm dV+c_0x^2\mathrm dU)+\mathrm dx^2+\mathrm dz^2+(\mathrm dw+f_0x\,\mathrm dU)^2\) solves the matrix equation exactly when \(c_0=-f_0^2/2\), that is \(\mathrm{Ric}^{(4)}_{ab}=\tfrac12F_{ac}F_b{}^c\) (exact). This fixes no \(G_{\mathrm{eff}}\), which needs an action or a loop. In~\cite{HoKawaiSteinacker2025}, the reading of matrices as momenta is described as not yet very successful for general relativity, the classical action is said not to be interpretable as the Einstein--Hilbert action, and that term is attributed to one loop on a noncommutative background; those statements concern the action and a vacuum without this fibre, and the result above is at the level of the equation of motion, so the two are not compared further.

These statements are for Equation~\eqref{eq:Smatrix} with \(C_{\mu\nu\lambda}=0\), tree level, the stated representations and first-order operators, the identities being exact at the sampled points and in the two relations. Nothing here derives that the matrix integral selects this class or its algebra \(\mathcal A\), or the frame and its non-degeneracy, which are properties of the class and are not derived from dynamics, so Assumptions~\ref{ass:frame-lock}, \ref{ass:nondeg} and~\ref{ass:lorentzian} are not discharged; Assumption~\ref{ass:local-lorentz}, a statement about the action~\eqref{eq:Smicro}, is not derived, since this class is not shown to reproduce Presentation~B. Domains of the unbounded operators, the regulated trace, the norm and masses of the torsion modes, higher-order and non-equivariant fluctuations, the bilocal modes of End of the Hilbert space, loops, supersymmetry and a value or sign of \(G_{\mathrm{eff}}\) are not treated, and the frame of this class is not identified with the emergent metric of Chapter~\ref{ch:vierbein-emergence}. Problems~D, F and~N remain open and PROGRAMME-CONDITIONAL.

